\documentclass[../DoAn.tex]{subfiles}
\begin{document}

\begin{center}
    \Large{\textbf{TÓM TẮT NỘI DUNG ĐỒ ÁN}}\\
\end{center}
\vspace{1cm}
Ô nhiễm không khí là mối đe dọa sức khỏe cộng đồng hàng đầu toàn cầu, đặc biệt nghiêm trọng tại Việt Nam khi nồng độ PM2.5 trung bình năm cao gấp gần sáu lần ngưỡng khuyến nghị của WHO. Hệ thống quan trắc quốc gia chỉ có khoảng 35 trạm đo tự động, không đủ phản ánh chất lượng không khí ở mức vi mô cho từng khu vực đô thị. Các giải pháp thương mại hiện có như PAM Air, IQAir và PurpleAir đều yêu cầu kết nối WiFi hoặc 4G tại mỗi điểm đo, chi phí thiết bị cao và phần mềm mã nguồn đóng. Không giải pháp nào đồng thời đáp ứng bốn tiêu chí: đo đa thông số, truyền dữ liệu tầm xa không cần WiFi, mã nguồn mở và chi phí thấp. Trên cơ sở đó, đồ án lựa chọn hướng tiếp cận xây dựng hệ thống IoT bốn tầng sử dụng cảm biến chi phí thấp kết hợp giao thức truyền thông LoRa tầm xa, loại bỏ yêu cầu hạ tầng mạng tại mỗi điểm đo.

Hệ thống gồm Sensor Node (vi điều khiển ESP32 kết hợp cảm biến PMS7003, CCS811, AHT10) đo sáu thông số môi trường, đóng gói thành gói tin nhị phân 18 byte và gửi qua LoRa đến Gateway. Gateway chuyển tiếp dữ liệu lên Backend Server (Node.js, TimescaleDB), nơi tính chỉ số AQI theo chuẩn US EPA và phát dữ liệu thời gian thực qua Socket.IO đến giao diện web React. Đồ án có bốn đóng góp chính: (i) hệ thống IoT đa tầng end-to-end hoạt động ổn định trong môi trường production với tổng quy mô 13.759 dòng mã, (ii) giao thức LoRa binary 18 byte giảm thời gian chiếm dụng kênh truyền khoảng 11 lần so với JSON kèm cơ chế chống trùng lặp dữ liệu, (iii) pipeline xử lý telemetry trên Backend đảm bảo không mất dữ liệu đo với toàn bộ 44 kịch bản kiểm thử đạt pass, và (iv) chi phí phần cứng khoảng 895.000 VNĐ mỗi Sensor Node — thấp hơn giải pháp IQAir khoảng 11 lần — với toàn bộ mã nguồn được công khai.
\begin{flushright}
Sinh viên thực hiện\\
\begin{tabular}{@{}c@{}}
\textit{(Ký và ghi rõ họ tên)}\\[1.5cm]
Nguyễn Đức Thắm
\end{tabular}
\end{flushright}

\end{document}